% Standalone proof. Shared notation and prose conventions: analyses/README.md.
\documentclass[11pt,a4paper]{article}
\usepackage[T1]{fontenc}
\usepackage[utf8]{inputenc}
\usepackage{amsmath,amsthm}
\usepackage{newtxtext,newtxmath}
\usepackage[a4paper,margin=21mm,headheight=15pt]{geometry}
\usepackage{microtype,booktabs,array,tabularx,longtable}
\usepackage[table]{xcolor}
\usepackage{tikz,fancyhdr,enumitem,seqsplit}
\definecolor{Teal}{HTML}{176B70}
\definecolor{Ink}{HTML}{203443}
\definecolor{Pale}{HTML}{EDF5F5}
\definecolor{Line}{HTML}{C7D8DB}
\definecolor{Muted}{HTML}{536974}
\definecolor{Warm}{HTML}{F8F4EC}
\usepackage[colorlinks=true,linkcolor=Teal,urlcolor=Teal]{hyperref}
\color{Ink}
\newcommand{\blank}{\lambda}
\newcommand{\digits}{\{0,1,\ldots,9\}}
\newcommand{\ph}[1]{\mathsf{#1}}
\newcommand{\EXPERT}{\mathcal E}
\newcommand{\CTRL}{\mathcal C}
\newcommand{\PTO}{\operatorname{PTO}}
\newcommand{\triple}[3]{(#1,#2,\mathsf{#3})}
\newcommand{\code}[1]{\texttt{#1}}
\newcommand{\hash}[1]{{\scriptsize\ttfamily\seqsplit{#1}}}
\newtheorem{theorem}{Theorem}
\newtheorem{lemma}{Lemma}
\newtheorem{corollary}{Corollary}
\setlength{\parindent}{0pt}
\setlength{\parskip}{6pt}
\setlist[itemize]{nosep,leftmargin=1.2em}
\pagestyle{fancy}
\fancyhf{}
\fancyhead[L]{\small\sffamily\color{Muted}Decimal Addition}
\fancyhead[R]{\small\sffamily\color{Muted}Metric Extraction}
\fancyfoot[L]{\scriptsize\sffamily\color{Muted}Exact PTO Equivalence}
\fancyfoot[R]{\small\thepage}
\renewcommand{\headrulewidth}{.35pt}
\renewcommand{\footrulewidth}{0pt}
\newcommand{\sectiontitle}[2]{\refstepcounter{section}\vspace*{2pt}{\LARGE\bfseries\color{Teal}\thesection\quad #2}\par\vspace{7pt}}
\date{}
\pdfinfoomitdate=1
\pdftrailerid{}
\begin{document}
{\Huge\bfseries Exact PTO Equivalence\\[3pt]Decimal Addition}\par
\vspace{10pt}
{\Large Metric-Extracted Controller}\par
{\small\color{Muted}\path{analyses/metric_extraction/DFST/add10/warmup_50000_steps_4000000/}}\par
\vspace{7pt}

\vspace{4pt}
\sectiontitle{1 \enspace STATEMENT AND SCOPE}{Statement and Scope}\label{sec:common1}
\textbf{Controllers.}
The exported table \path{minimized_transition_table.pdf} defines controller
$\CTRL$, with initial state $0$ and halt state $51$.
The function \code{add\_grid(10)} defines expert $\EXPERT$, with initial
state $0$ and halt state $18$. The adjacent expert export matches
this function entry for entry. All tables retain the exported state numbers.

\textbf{Input Domain.}
Let $D=\digits$. Row $0$ initially contains the contiguous word $a+b$,
starting at column $0$, with $a\in D^m$, $b\in D^n$, and $m,n\ge1$.
All remaining cells contain blanks. The head starts at $(0,0)$.
Write $A=\operatorname{val}_{10}(a)$ and $B=\operatorname{val}_{10}(b)$.
The domain permits arbitrary operand lengths, unequal lengths, and leading zeros.

\textbf{Execution Convention.}
The grid uses $\mathbb Z^2$, with coordinates increasing rightward and
downward. The symbol $\blank$ denotes a blank. A transition
$(q',w,M)$ writes $w$, performs motion $M$, and enters state $q'$.
Motions $L,R,D,U$ move one cell. Motion $S$ writes and stops without moving.

\textbf{Policy-Trajectory Observations.}
For input $a+b$, define the policy-trajectory observation (PTO) by
\[
\PTO(a+b)=\bigl((r_t,w_t,M_t)\bigr)_{t=0}^{T-1},
\]
where $r_t$, $w_t$, and $M_t$ denote the read symbol, written symbol,
and motion at controller step $t$. A finite execution contributes the final
$S$ action exactly once, without halt-loop padding.

\begin{theorem}[Exact PTO Equivalence]\label{thm:main}
For every input in the stated domain,
$\PTO_{\CTRL}(a+b)=\PTO_{\EXPERT}(a+b)$.
Both executions terminate on the same step. The complete grids and head
positions agree at every step, including the final configuration.
\end{theorem}

The theorem concerns autonomous executions of the exported discrete controller
on the stated input domain.
\clearpage
\sectiontitle{2 \enspace EXPERT BEHAVIOR}{Expert Behavior}\label{sec:common2}
\subsection{Accumulator and Marker Invariant}\label{sec:part1}
Write $B=\sum_{j=0}^{n-1}b_j10^j$, with $b_0$ the rightmost digit.
The expert scans to the separator at column $m$, steps left and down,
writes a working $+$ below column $m-1$, and enters state $16$ above the marker.

\begin{lemma}[Round Invariant]\label{lem:round}
After processing $k$ right-operand digits, $0\le k\le n$, the expert enters
state $16$ at $(p_k,0)$, where $p_k=m-1-k$. At this point:
\begin{itemize}
\item Row $1$ contains exactly one nonblank symbol: a $+$ at column $p_k$.
\item Row $0$ contains a contiguous accumulator on $[\ell_k,m-1]$, with value
$C_k=A+\sum_{j<k}b_j10^j$ and left edge $\ell_k\le p_k+1$.
Column $m-1$ carries the units digit.
\item The separator remains at column $m$. The unprocessed prefix of $b$
occupies columns $m+1,\ldots,m+n-k$. All remaining cells contain blanks.
\end{itemize}
\end{lemma}
\begin{proof}
Initialization gives $k=0$, $\ell_0=0$, $C_0=A$, and $p_0=m-1$.
State $16$ preserves the digit at $p_k$, or writes $0$ on a blank.
A blank at $p_k$ forces $\ell_k=p_k+1$, so zero extension preserves contiguity.

For $k<n$, the right scan reaches the blank after the unprocessed prefix,
steps left to digit $b_k$ at $r_k=m+n-k$, erases $b_k$, and descends.
The expert transports $b_k$ left along row $1$ to the marker. The distance satisfies
\[
r_k-p_k=n+1\ge2. \tag{1}
\]
Consequently, the first two transport steps read blanks. At the marker, the
expert ascends and adds $b_k$ to the accumulator digit $e$ at column $p_k$.
If $e+b_k<10$, the expert writes $e+b_k$ and descends. Otherwise, the expert
writes $e+b_k-10$, moves left, and propagates one carry. Each digit $9$
changes to $0$; the first smaller digit increases by one, or a blank changes
to $1$. The update adds exactly $b_k10^k$ and preserves contiguity.

The expert returns along row $1$, erases the marker, steps left, writes the
next marker at $p_{k+1}=p_k-1$, and ascends into state $16$.
The earlier zero extension gives $\ell_{k+1}\le p_k=p_{k+1}+1$.
The accumulator now has value $C_{k+1}$, establishing the next round.
\end{proof}

\begin{lemma}[Termination and Final Configuration]\label{lem:finish}
The expert terminates with only the padded base-$10$ representation of $A+B$
on row $0$, ending at column $m-1$. The final head stops at $(m-n-1,0)$.

Writing $\ell_{10}(s)$ for the decimal digit count, with $\ell_{10}(0)=1$,
the output width equals $L=\max\{m,n+1,\ell_{10}(A+B)\}$.
\end{lemma}
\begin{proof}
Each round consumes one of the $n$ right-operand digits. Every scan crosses
a finite interval, and every carry reaches a digit below $9$ or a blank.
After round $n$, state $16$ ensures a digit at $p_n$. The next right scan
returns to the separator. State $3$ erases the separator and descends.
State $17$ walks left, erases the final marker, ascends, preserves the
accumulator digit, and stops. The invariant gives $C_n=A+B$, and cleanup
removes all work symbols.

The original $m$ columns persist, and marker extension reaches $p_n$,
ensuring at least $n+1$ digits. Every additional carry extension contributes
a digit required by the current numerical value. Since partial sums never
decrease, the final width equals $L$.
\end{proof}

\clearpage
\sectiontitle{3 \enspace OBSERVATION PROTOCOL}{Observation Protocol}\label{sec:common3}
The observation protocol refines expert states by geometric context. Each phase records the permitted read symbols and successor phases. Section~\ref{sec:common4} gives the complete protocol together with the controller-state relation; the shared tables avoid repeating the same phase definitions.

\begin{lemma}[Protocol Coverage]\label{lem:coverage}
Every expert execution on a valid decimal addition grid follows the $59$
phases in Section~\ref{sec:common4}, starting at phase $0$ and ending at phase $58$.
\end{lemma}
\begin{proof}
Initialization satisfies the first phases because both operands contain a digit.
Lemma~\ref{lem:round} gives the contiguous accumulator, the ensured digit,
the single working marker and the remaining right-operand prefix in every
round. The explanations in Sections~\ref{sec:part2} and~\ref{sec:part3} establish all first-cell,
boundary, transport and return readings. The carry rule processes exactly
the digits permitted in phase $50$. Induction over the transitions of each
round preserves these conditions. Lemma~\ref{lem:finish} establishes finite
cleanup and the final phase.
\end{proof}

\clearpage
\sectiontitle{4 \enspace SIMULATION CERTIFICATE}{Simulation Certificate}\label{sec:common4}
\subsection{Scanning, Carrying, and Cleanup}\label{sec:part2}
A proof phase records geometric context. Phase $i$ pairs expert state
$g_i$ with controller state $q\in Q_i$, with identical grids and heads. An entry $x:(j,w,M)$ means read $x$, write $w$, move $M$, and
advance to phase $j$. The index $j$ denotes a phase. Each $Q_i$ lists the related controller states.

Here $e\in D$ and $h\in\{0,\ldots,8\}$. A digit denotes the corresponding symbol when read or written. Phases $10$--$49$ comprise ten digit-transport families in the next certificate subsection, with $F_d=10+4d$.

\begin{center}
\fontsize{9.1}{11}\selectfont\renewcommand{\arraystretch}{1.3}
\setlength{\tabcolsep}{4pt}
\begin{tabularx}{\textwidth}{r p{4.8cm} c c X}
\toprule\rowcolor{Pale}$i$ & Phase & $g_i$ & $Q_i$ & Read $:$ (next phase, write, move)\\\midrule
0 & First left-operand digit & 0 & $\{0\}$ & $e:\triple{1}{e}{R}$\\
1 & Find the input separator & 0 & $\{0\}$ & $e:\triple{1}{e}{R}\quad +:\triple{2}{+}{L}$\\
2 & Descend below units digit & 1 & $\{1\}$ & $e:\triple{3}{e}{D}$\\
3 & Place the working marker & 1 & $\{2\}$ & $\blank:\triple{4}{+}{U}$\\
4 & Ensure a digit above marker & 16 & $\{3\}$ & $e:\triple{6}{e}{R}\quad\blank:\triple{5}{0}{R}$\\
5 & Adjacent existing digit & 2 & $\{4\}$ & $e:\triple{6}{e}{R}$\\
6 & Scan accumulator to separator & 2 & $\{4\}$ & $e:\triple{6}{e}{R}\quad +:\triple{7}{+}{R}$\\
7 & Inspect remaining right operand & 2 & $\{5\}$ & $e:\triple{8}{e}{R}\quad\blank:\triple{54}{\blank}{L}$\\
8 & Scan remaining right operand & 2 & $\{5\}$ & $e:\triple{8}{e}{R}\quad\blank:\triple{9}{\blank}{L}$\\
9 & Read rightmost factor digit & 3 & $\{6\}$ & $d:\triple{F_d}{\blank}{D}\quad(d\in D)$\\
\midrule
50 & Propagate decimal carry & 14 & $\{11\}$ & $9:\triple{50}{0}{L}\quad h:\triple{51}{h+1}{D}$\newline $\blank:\triple{51}{1}{D}$\\
51 & First return cell after carry & 15 & $\{15\}$ & $\blank:\triple{52}{\blank}{R}$\\
52 & Return to working marker & 15 & $\{15\}$ & $\blank:\triple{52}{\blank}{R}\quad +:\triple{3}{\blank}{L}$\\
53 & Erase and shift marker & 15 & $\{11,15\}$ & $+:\triple{3}{\blank}{L}$\\
54 & Erase input separator & 3 & $\{6\}$ & $+:\triple{55}{\blank}{D}$\\
55 & Begin cleanup below separator & 17 & $\{24\}$ & $\blank:\triple{56}{\blank}{L}$\\
56 & Erase final working marker & 17 & $\{25\}$ & $\blank:\triple{56}{\blank}{L}\quad +:\triple{57}{\blank}{U}$\\
57 & Preserve final digit and stop & 17 & $\{26\}$ & $e:\triple{58}{e}{S}$\\
58 & Halted & 18 & $\{51\}$ & --\\
\bottomrule
\end{tabularx}
\end{center}

\textbf{Why These Read Symbols Cover Every Valid Run.}
Both operands contain at least one digit. Initialization therefore reads a digit before
the separator, descends below a digit and places the marker on a blank.
Lemma~\ref{lem:round} places an existing digit immediately after each zero extension. The scan crosses the contiguous accumulator to the separator.
Immediately to the separator's right, a blank means the next left step
reads that separator; after any digit, the next left step instead reads a digit.

A carry moves strictly left of the marker, so the first return cell on
row $1$ contains a blank. A noncarry addition descends directly onto the marker.
Erasing and shifting the marker reaches a blank. Cleanup
starts below the separator on a blank, walks to the marker, then ascends
to a digit. These facts establish the table's restrictions.

\newpage
\subsection{Digit-Transport Phases}\label{sec:part3}
For each $d\in D$, four phases carry the selected digit to the accumulator:
\[
F_d=10+4d,\qquad H_d=F_d+1,\qquad W_d=F_d+2,\qquad A_d=F_d+3.
\]
All four phases use expert state $4+d$, paired with singleton controller
sets containing $f_d,h_d,w_d,a_d$, respectively. The protocol follows:

\begin{center}
\small\renewcommand{\arraystretch}{1.35}
\begin{tabular}{lll}
\toprule\rowcolor{Pale}Phase & Meaning & Read $:$ (next phase, write, move)\\\midrule
$F_d$ & First transport blank & $\blank:\triple{H_d}{\blank}{L}$\\
$H_d$ & Second transport blank & $\blank:\triple{W_d}{\blank}{L}$\\
$W_d$ & Walk left to the marker & $\blank:\triple{W_d}{\blank}{L}\quad +:\triple{A_d}{+}{U}$\\
$A_d$ & Add at the marked column & $e:\triple{53}{d+e}{D}$ if $d+e<10$\\
 & & $e:\triple{50}{d+e-10}{L}$ if $d+e\ge10$\\
\bottomrule
\end{tabular}
\end{center}

\textbf{The Actual Controller-State Numbers.}
Each row retains the state numbers from the exported controller table. This table expands the entire $40$-phase family.

\begin{center}
\small\renewcommand{\arraystretch}{1.3}\setlength{\tabcolsep}{12pt}
\begin{tabular}{rrrrrr}
\toprule\rowcolor{Pale}$d$ & Expert $4+d$ & $f_d$ & $h_d$ & $w_d$ & $a_d$\\\midrule
0 & 4 & 27 & 28 & 29 & 30 \\
1 & 5 & 12 & 13 & 14 & 11 \\
2 & 6 & 39 & 40 & 41 & 42 \\
3 & 7 & 31 & 32 & 33 & 34 \\
4 & 8 & 43 & 44 & 45 & 46 \\
5 & 9 & 47 & 48 & 49 & 50 \\
6 & 10 & 7 & 8 & 9 & 10 \\
7 & 11 & 20 & 21 & 22 & 23 \\
8 & 12 & 16 & 17 & 18 & 19 \\
9 & 13 & 35 & 36 & 37 & 38 \\
\bottomrule
\end{tabular}
\end{center}

\textbf{Transport Coverage.}
Equation~(1) gives $r_k-p_k=n+1\ge2$. The first two transport cells therefore contain blanks. After two left
steps, the head reaches either the marker or a blank to the marker's right. The walk reaches the
marker and ascends to a digit ensured at the start of the round.
The two branches implement decimal addition with at most one carry:
\[
d+e=s+10c,\qquad s\in D,\quad c\in\{0,1\}.
\]
This accounts for every pair $(d,e)$, including $0$, $9$ and all carry cases.

\textbf{Multiple Roles of Controller State $11$.}
State $11$ performs addition of $1$ when reading a digit, propagates carries when
entered to the left of an overflowing column, and erases the marker when
reading $+$. With $d=1$ and no carry, the controller enters state $15$; every other noncarry branch enters state $11$. Both belong to
$Q_{53}=\{11,15\}$ and have the same required action on $+$.
Every overflow enters state $11\in Q_{50}$.
The phase relation distinguishes these geometric contexts and therefore
does not require a one-to-one correspondence between controller states.

\textbf{Local Transition Identity.}
For every phase $i$, controller state $q\in Q_i$ and listed branch
$x:(j,w,M)$, direct comparison with the exported tables gives
\begin{equation}
\delta_\EXPERT(g_i,x)=(g_j,w,M),\qquad
\delta_\CTRL(q,x)=(q',w,M),\quad q'\in Q_j. \tag{2}
\end{equation}
Every required entry exists. A transition reaches a halt state exactly
when $j=58$, with $M=S$. The certificate contains $60$ phase--state pairs,
including the terminal pair, and $266$ local checks of (2). The checks cover
all $223$ defined nonterminal controller entries. The remaining $12$ entries
contain halt loops, excluded from PTOs.

\clearpage
\sectiontitle{5 \enspace EQUIVALENCE PROOF}{Equivalence Proof}\label{sec:common5}
\label{sec:equivalence}
\begin{proof}[Proof of Theorem~\ref{thm:main}]
Initially both configurations have the same grid and head,
with control states $0$ related by phase $0$. Suppose the relation holds
before a controller step. Lemma~\ref{lem:coverage} supplies an allowed read symbol
$x$. Equal grids and heads give the same read symbol to both controllers.
Equation~(2) gives equal writes and motions and preserves the relation
between successor states. The next grids and heads therefore agree.

Induction establishes equality of every finite execution prefix. The terminal
condition permits a halt exactly on an edge to the terminal phase.
Lemma~\ref{lem:finish} establishes termination.
Both controllers therefore halt together with equal complete PTOs and final
configurations.
\end{proof}

Protocol coverage and the local transition identity exclude every undefined entry on the stated input domain.
\begin{center}
\small\renewcommand{\arraystretch}{1.3}
\begin{tabular}{llll}
\toprule\rowcolor{Pale}Input & Final digit string & Occupied columns & Final head\\\midrule
$0+0$ & $00$ & $[-1,0]$ & $(-1,0)$\\
$9+9$ & $18$ & $[-1,0]$ & $(-1,0)$\\
$999+1$ & $1000$ & $[-1,2]$ & $(1,0)$\\
$0+999$ & $0999$ & $[-3,0]$ & $(-3,0)$\\
$001+09$ & $010$ & $[0,2]$ & $(0,0)$\\
\bottomrule
\end{tabular}
\end{center}

\clearpage
\sectiontitle{6 \enspace VERIFICATION AND ARTIFACTS}{Verification and Artifacts}\label{sec:common6}
The universal claim follows from the geometric invariant, protocol coverage
and induction. The following checks bind the finite certificate to the exact
adjacent files and independently exercise the controller.

\textbf{Finite Certificate.}
Run
\begin{center}\code{python3 verify\_equivalence.py {-}{-}quick}\end{center}
from this proof directory. This uses only the Python standard library and checks:
\begin{itemize}
\item SHA-256 fingerprints of the controller exports, the expert-source
snapshots and the finite protocol;
\item exact agreement between the minimized CSV and JSON controllers;
\item exact agreement between the exported expert and \code{add\_grid(10)};
\item all $266$ local obligations for the $60$ related phase--state pairs,
including start, halt and successor closure;
\item coverage of all $223$ defined nonterminal controller entries.
\end{itemize}
The file \path{equivalence_verification.json} records the phase relation
and every checked transition.

\textbf{Independent Complete Executions.}
Omitting \code{{-}{-}quick} also runs two separate grid memories and compares
every read symbol, write, move and head position. The checker compares the complete final grids against the exact padded decimal sum and final-head formula.
The validation set contains all $12{,}100$ pairs of operand strings of
lengths $1$--$2$, plus $44$ directed cases and $256$ fixed-seed random cases.
Directed cases include long all-$9$ carry chains, zeros, unequal lengths,
leading zeros and operands of length $1024$.

All $3{,}577{,}109$ compared PTO steps agree. The execution checks exercise every defined nonterminal controller entry. The finite tests supplement the length-independent proof.

\textbf{Printed-Table Verification.}
Printed-table verification compares all $624$ cells of the two-page PDF
against the CSV: $235$ defined entries and $389$ undefined cells.
The optional command
\begin{center}\code{python3 verify\_equivalence.py {-}{-}quick {-}{-}check-pdf}\end{center}
repeats this printed-table check with \code{pypdf} installed.

\textbf{Artifact Identity.}
The title page gives the repository-relative case directory. This proof directory
contains the controller and expert exports, reference code, protocol, checker,
and saved verification report. The verification script stores or reads SHA-256 checksums for these artifacts. A checksum identifies the exact bytes
of a file; the checker compares the stored and computed values to detect changes.
These checks bind the certificate to the supplied exports. The invariants,
protocol coverage, and local transition identity establish equivalence.
\end{document}
